\documentclass[UTF8]{ctexart}
\usepackage[a4paper,margin=1.8cm]{geometry}
\usepackage{graphicx}
\usepackage{float}
\usepackage{longtable}
\usepackage{array}
\usepackage{xcolor}
\usepackage{fancyvrb}
\usepackage{hyperref}
\setlength{\parindent}{0pt}
\setlength{\parskip}{0.45em}
\begin{document}
\section*{补充模拟卷 01：自顶向下分析与 LR 概念补缺}

\subsection*{A 卷题目}

\subsubsection*{一、递归下降、回溯与左公因子（20 分）}

文法：

\begin{Verbatim}[fontsize=\small,breaklines=true]
S -> a A | a B
A -> c
B -> d
\end{Verbatim}

1. 说明该文法在递归下降分析中为什么可能需要回溯。 2. 提取左公因子。 3. 判断提取左公因子后的文法是否适合构造 LL(1) 预测分析表。

\subsubsection*{二、二义性文法与预测分析表（20 分）}

文法：

\begin{Verbatim}[fontsize=\small,breaklines=true]
E -> E + E | E * E | id
\end{Verbatim}

1. 说明该文法为什么不适合直接构造 LL(1) 预测分析表。 2. 指出它至少存在的两个问题。 3. 简述通常需要怎样处理。

\subsubsection*{三、句柄与归约序列（20 分）}

文法：

\begin{Verbatim}[fontsize=\small,breaklines=true]
S -> A B
A -> a
B -> b
\end{Verbatim}

对输入串：

\begin{Verbatim}[fontsize=\small,breaklines=true]
a b
\end{Verbatim}

1. 写出自底向上的归约序列。 2. 说明该归约序列与最右推导的关系。

\subsubsection*{四、规约-规约冲突（20 分）}

某 LR(0) 状态中有：

\begin{Verbatim}[fontsize=\small,breaklines=true]
A -> id .
B -> id .
\end{Verbatim}

1. 说明为什么这是规约-规约冲突。 2. 若 \texttt{FOLLOW(A) = \textbackslash{}\{ ; \textbackslash{}\}}，\texttt{FOLLOW(B) = \textbackslash{}\{ ) \textbackslash{}\}}，说明 SLR 如何处理。 3. 若 \texttt{FOLLOW(A)} 和 \texttt{FOLLOW(B)} 有公共符号，会发生什么。

\subsubsection*{五、内核项、非内核项与 SLR 弱点（20 分）}

回答：

1. 什么是 LR(0) 项集中的内核项。 2. 什么是非内核项。 3. SLR 为什么比 LR(0) 强。 4. SLR 为什么仍然有弱点。

\subsection*{B 按考试标准重写答案}

\subsubsection*{一、递归下降、回溯与左公因子答案}

\paragraph*{1. 为什么可能需要回溯}

原文法中 \texttt{S} 有两条产生式：

\begin{Verbatim}[fontsize=\small,breaklines=true]
S -> a A
S -> a B
\end{Verbatim}

递归下降分析在选择产生式时，通常先看当前输入符号。若输入串以 \texttt{a} 开头，则两条产生式的第一个终结符都是 \texttt{a}：

\begin{longtable}{|p{0.47\linewidth}|p{0.47\linewidth}|}
\hline
产生式 & 可见的第一个输入符号 \\
\hline
\endfirsthead
产生式 & 可见的第一个输入符号 \\
\hline
\endhead
\texttt{S -> a A} & \texttt{a} \\
\hline
\texttt{S -> a B} & \texttt{a} \\
\hline
\end{longtable}

因此只看 1 个输入符号时无法确定应选择哪条产生式。若分析器先选 \texttt{S -> a A}，但后续输入是 \texttt{a d}，则 \texttt{A -> c} 无法匹配 \texttt{d}，必须退回去改选 \texttt{S -> a B}。这就是回溯来源。

\paragraph*{2. 提取左公因子}

把公共前缀 \texttt{a} 提出来：

\begin{Verbatim}[fontsize=\small,breaklines=true]
S  -> a S'
S' -> A | B
A  -> c
B  -> d
\end{Verbatim}

由于 \texttt{A -> c}、\texttt{B -> d}，也可以等价写成：

\begin{Verbatim}[fontsize=\small,breaklines=true]
S  -> a S'
S' -> c | d
\end{Verbatim}

\paragraph*{3. LL(1) 判断}

对提取左公因子后的文法计算 FIRST：

\begin{longtable}{|p{0.47\linewidth}|p{0.47\linewidth}|}
\hline
非终结符或右部 & FIRST \\
\hline
\endfirsthead
非终结符或右部 & FIRST \\
\hline
\endhead
\texttt{c} & \texttt{\textbackslash{}\{ c \textbackslash{}\}} \\
\hline
\texttt{d} & \texttt{\textbackslash{}\{ d \textbackslash{}\}} \\
\hline
\texttt{S' -> c} & \texttt{\textbackslash{}\{ c \textbackslash{}\}} \\
\hline
\texttt{S' -> d} & \texttt{\textbackslash{}\{ d \textbackslash{}\}} \\
\hline
\texttt{S -> a S'} & \texttt{\textbackslash{}\{ a \textbackslash{}\}} \\
\hline
\end{longtable}

\texttt{S'} 的两条产生式 FIRST 集不相交：

\begin{Verbatim}[fontsize=\small,breaklines=true]
FIRST(c) intersection FIRST(d) = empty
\end{Verbatim}

因此预测分析表中不会在同一表项填入两条产生式。对应表项为：

\begin{longtable}{|p{0.19\linewidth}|p{0.19\linewidth}|p{0.19\linewidth}|p{0.19\linewidth}|p{0.19\linewidth}|}
\hline
非终结符 & 输入 a & 输入 c & 输入 d & 输入 \$ \\
\hline
\endfirsthead
非终结符 & 输入 a & 输入 c & 输入 d & 输入 \$ \\
\hline
\endhead
S & \texttt{S -> a S'} & error & error & error \\
\hline
S' & error & \texttt{S' -> c} & \texttt{S' -> d} & error \\
\hline
\end{longtable}

所以提取左公因子后的文法适合构造 LL(1) 预测分析表。

\subsubsection*{二、二义性文法与预测分析表答案}

\paragraph*{1. 为什么不能直接构造 LL(1) 表}

文法：

\begin{Verbatim}[fontsize=\small,breaklines=true]
E -> E + E | E * E | id
\end{Verbatim}

不能直接构造 LL(1) 预测分析表，原因有两个层次。

第一，它有直接左递归：

\begin{Verbatim}[fontsize=\small,breaklines=true]
E -> E + E
E -> E * E
\end{Verbatim}

递归下降分析 \texttt{E} 时会立即再次调用 \texttt{E}，可能无限递归。

第二，它没有规定 \texttt{+}、\texttt{*} 的优先级和结合性，因此是二义性的。例如输入：

\begin{Verbatim}[fontsize=\small,breaklines=true]
id + id * id
\end{Verbatim}

可以有两种结构：

\begin{longtable}{|p{0.47\linewidth}|p{0.47\linewidth}|}
\hline
分析结构 & 含义 \\
\hline
\endfirsthead
分析结构 & 含义 \\
\hline
\endhead
\texttt{(id + id) * id} & 先算加法，再算乘法 \\
\hline
\texttt{id + (id * id)} & 先算乘法，再算加法 \\
\hline
\end{longtable}

同一个输入串有不同语法树，所以不适合直接做确定的预测分析。

\paragraph*{2. 至少存在的两个问题}

\begin{longtable}{|p{0.31\linewidth}|p{0.31\linewidth}|p{0.31\linewidth}|}
\hline
问题 & 具体表现 & 对 LL(1) 的影响 \\
\hline
\endfirsthead
问题 & 具体表现 & 对 LL(1) 的影响 \\
\hline
\endhead
直接左递归 & \texttt{E -> E + E}、\texttt{E -> E * E} & 递归下降会无限递归 \\
\hline
二义性 & 没有体现优先级和结合性 & 同一个输入串有多种语法树 \\
\hline
FIRST 冲突 & 三条产生式都可能以 \texttt{id} 开始 & 预测分析表的同一格会出现多条产生式 \\
\hline
\end{longtable}

考试中写出前两个问题即可完整回答第 2 问；写出 FIRST 冲突可以补强说明。

\paragraph*{3. 通常处理方式}

先改写成表达优先级和结合性的非二义文法，再消除左递归。常用改写为：

\begin{Verbatim}[fontsize=\small,breaklines=true]
E -> E + T | T
T -> T * F | F
F -> id
\end{Verbatim}

该文法体现了 \texttt{*} 的优先级高于 \texttt{+}。为了构造 LL(1) 表，再消除左递归：

\begin{Verbatim}[fontsize=\small,breaklines=true]
E  -> T E'
E' -> + T E' | epsilon
T  -> F T'
T' -> * F T' | epsilon
F  -> id
\end{Verbatim}

最后再计算 FIRST、FOLLOW 并填写预测分析表。

\subsubsection*{三、句柄与归约序列答案}

\paragraph*{1. 自底向上的归约序列}

文法：

\begin{Verbatim}[fontsize=\small,breaklines=true]
S -> A B
A -> a
B -> b
\end{Verbatim}

输入串是：

\begin{Verbatim}[fontsize=\small,breaklines=true]
a b
\end{Verbatim}

自底向上归约时，每一步选择当前句型中的句柄，即某个产生式右部，然后把它归约为产生式左部：

\begin{longtable}{|p{0.23\linewidth}|p{0.23\linewidth}|p{0.23\linewidth}|p{0.23\linewidth}|}
\hline
当前句型 & 句柄 & 使用产生式 & 归约后 \\
\hline
\endfirsthead
当前句型 & 句柄 & 使用产生式 & 归约后 \\
\hline
\endhead
\texttt{a b} & \texttt{a} & \texttt{A -> a} & \texttt{A b} \\
\hline
\texttt{A b} & \texttt{b} & \texttt{B -> b} & \texttt{A B} \\
\hline
\texttt{A B} & \texttt{A B} & \texttt{S -> A B} & \texttt{S} \\
\hline
\end{longtable}

所以归约序列为：

\begin{Verbatim}[fontsize=\small,breaklines=true]
a b
=> A b
=> A B
=> S
\end{Verbatim}

\paragraph*{2. 与最右推导的关系}

自底向上归约是最右推导的逆过程。上面归约序列反过来就是一个最右推导：

\begin{Verbatim}[fontsize=\small,breaklines=true]
S
=> A B
=> A b
=> a b
\end{Verbatim}

检查“最右”的依据是：每次推导都替换当前句型中最右边的非终结符。

\begin{longtable}{|p{0.47\linewidth}|p{0.47\linewidth}|}
\hline
推导步骤 & 被替换的最右非终结符 \\
\hline
\endfirsthead
推导步骤 & 被替换的最右非终结符 \\
\hline
\endhead
\texttt{S => A B} & \texttt{S} \\
\hline
\texttt{A B => A b} & \texttt{B} \\
\hline
\texttt{A b => a b} & \texttt{A} \\
\hline
\end{longtable}

因此题目中的归约序列确实是该最右推导的逆序。

\subsubsection*{四、规约-规约冲突答案}

\paragraph*{1. 为什么是规约-规约冲突}

某 LR(0) 状态中同时有两个完成项：

\begin{Verbatim}[fontsize=\small,breaklines=true]
A -> id .
B -> id .
\end{Verbatim}

点在产生式末尾表示已经识别出完整右部，可以规约。在 LR(0) 中，完成项不看向前输入符号，会尝试在所有可能输入上填规约动作。因此同一个状态会同时要求：

\begin{Verbatim}[fontsize=\small,breaklines=true]
reduce A -> id
reduce B -> id
\end{Verbatim}

同一个 ACTION 表项不能同时放两个规约动作，所以这是规约-规约冲突。

\paragraph*{2. SLR 如何处理给定 FOLLOW 集}

SLR 的规则是：完成项 \texttt{A -> alpha .} 只在 \texttt{FOLLOW(A)} 中的终结符上填规约。

题目给出：

\begin{Verbatim}[fontsize=\small,breaklines=true]
FOLLOW(A) = { ; }
FOLLOW(B) = { ) }
\end{Verbatim}

因此填表为：

\begin{longtable}{|p{0.47\linewidth}|p{0.47\linewidth}|}
\hline
输入符号 & ACTION \\
\hline
\endfirsthead
输入符号 & ACTION \\
\hline
\endhead
\texttt{;} & \texttt{r(A -> id)} \\
\hline
\texttt{)} & \texttt{r(B -> id)} \\
\hline
\end{longtable}

因为 \texttt{\textbackslash{}\{ ; \textbackslash{}\}} 与 \texttt{\textbackslash{}\{ ) \textbackslash{}\}} 不相交，所以这两个规约动作不会落在同一个 ACTION 表项中，SLR 可以消除这个 LR(0) 状态中的冲突。

\paragraph*{3. FOLLOW 有公共符号时的结果}

若存在某个终结符 \texttt{x}，使得：

\begin{Verbatim}[fontsize=\small,breaklines=true]
x in FOLLOW(A)
x in FOLLOW(B)
\end{Verbatim}

则 SLR 会同时填写：

\begin{Verbatim}[fontsize=\small,breaklines=true]
ACTION[状态, x] = r(A -> id)
ACTION[状态, x] = r(B -> id)
\end{Verbatim}

同一表项仍有两个规约动作，规约-规约冲突仍然存在。

\subsubsection*{五、内核项、非内核项与 SLR 弱点答案}

\paragraph*{1. 内核项}

LR(0) 项集中的内核项包括两类：

\begin{Verbatim}[fontsize=\small,breaklines=true]
1. 增广文法的初始项 S' -> . S。
2. 所有点不在产生式最左端的 LR(0) 项。
\end{Verbatim}

例子：

\begin{longtable}{|p{0.31\linewidth}|p{0.31\linewidth}|p{0.31\linewidth}|}
\hline
项 & 是否内核项 & 原因 \\
\hline
\endfirsthead
项 & 是否内核项 & 原因 \\
\hline
\endhead
\texttt{S' -> . S} & 是 & 初始项特例 \\
\hline
\texttt{A -> a . B} & 是 & 点不在最左端 \\
\hline
\texttt{A -> a B .} & 是 & 点不在最左端 \\
\hline
\end{longtable}

\paragraph*{2. 非内核项}

非内核项是除初始项以外，点在产生式最左端的项：

\begin{longtable}{|p{0.31\linewidth}|p{0.31\linewidth}|p{0.31\linewidth}|}
\hline
项 & 是否非内核项 & 原因 \\
\hline
\endfirsthead
项 & 是否非内核项 & 原因 \\
\hline
\endhead
\texttt{A -> . a B} & 是 & 点在最左端，且不是 \texttt{S' -> . S} \\
\hline
\texttt{B -> . c} & 是 & 点在最左端 \\
\hline
\end{longtable}

这些项通常由 closure 运算自动补入。

\paragraph*{3. SLR 为什么比 LR(0) 强}

LR(0) 看到完成项 \texttt{A -> alpha .} 后，不使用向前输入符号，规约范围过大。SLR 增加 FOLLOW 集限制：

\begin{Verbatim}[fontsize=\small,breaklines=true]
LR(0): A -> alpha . 在所有终结符上规约
SLR : A -> alpha . 只在 FOLLOW(A) 上规约
\end{Verbatim}

因此 SLR 可以消除一部分 LR(0) 中的移进-规约冲突和规约-规约冲突。

\paragraph*{4. SLR 为什么仍有弱点}

SLR 的 FOLLOW 集只与非终结符有关，不区分项目所在的具体 LR 状态和更精确的调用上下文。若 \texttt{FOLLOW(A)} 在全局范围内很大，SLR 会在某些当前状态不该规约的位置也填入规约动作。

对比：

\begin{longtable}{|p{0.31\linewidth}|p{0.31\linewidth}|p{0.31\linewidth}|}
\hline
方法 & 规约依据 & 精确程度 \\
\hline
\endfirsthead
方法 & 规约依据 & 精确程度 \\
\hline
\endhead
LR(0) & 不看输入符号 & 最粗 \\
\hline
SLR & 看 FOLLOW(A) & 比 LR(0) 精细，但 FOLLOW 是全局集合 \\
\hline
LR(1) & 看项目自身携带的向前看符号 & 更精确 \\
\hline
\end{longtable}

因此有些文法不是 SLR(1)，但可以是 LR(1) 或 LALR(1)。
\end{document}
